\chapterbackmatter{Solutions to Weinberg Exercises}
\ExerciseEditorialNote

\begin{enumerate}
  \WeinbergSolution{1}
  Let
  \[
    H=\begin{pmatrix}1&0\\0&-1\end{pmatrix},
    \qquad
    I=\begin{pmatrix}1&0\\0&1\end{pmatrix},
    \qquad
    F_+=\begin{pmatrix}0&1\\0&0\end{pmatrix},
    \qquad
    F_-=\begin{pmatrix}0&0\\1&0\end{pmatrix}.
  \]
  Declare the diagonal matrices \(H,I\) even and the off-diagonal matrices
  \(F_+,F_-\) odd.  The graded bracket is the commutator unless both
  entries are odd, in which case it is the anticommutator.  Direct
  multiplication gives
  \[
    [H,F_+]=2F_+,\qquad [H,F_-]=-2F_-,
    \qquad \{F_+,F_-\}=I,
  \]
  together with
  \[
    [H,I]=[I,F_\pm]=0,\qquad
    \{F_+,F_+\}=\{F_-,F_-\}=0.
  \]
  Every bracket is therefore again in the span of the four matrices.  This
  is the matrix Lie superalgebra \(\mathfrak{gl}(1|1)\), written in a
  homogeneous basis; associativity of matrix multiplication guarantees the
  graded Jacobi identity.

  \WeinbergSolution{2}
  We use real two-component Majorana spinors in \(2+1\) dimensions.  The
  assumptions of the Coleman--Mandula theorem first restrict every even
  generator, apart from the Poincar\'e generators, to be a Lorentz scalar
  internal generator.  If an odd generator transformed in Lorentz spin
  \(A>1/2\), its anticommutator with its adjoint would contain an even
  component of spin \(2A>1\).  Such a conserved tensor generator is
  forbidden.  A non-scalar odd generator must consequently have \(A=1/2\);
  call the resulting Majorana charges \(Q^I_\alpha\), where
  \(\alpha=1,2\) and \(I=1,\ldots,N\).

  Lorentz covariance, symmetry under
  \((\alpha,I)\leftrightarrow(\beta,J)\), and positivity then give, after
  an orthogonal change of basis and a conventional normalization,
  \[
    \{Q^I_\alpha,Q^J_\beta\}
      =2\delta^{IJ}(\gamma^\mu C)_{\alpha\beta}P_\mu
       +2C_{\alpha\beta}Z^{IJ},
    \qquad Z^{IJ}=-Z^{JI}.
  \]
  Here \(C_{\alpha\beta}\) is antisymmetric, whereas
  \((\gamma^\mu C)_{\alpha\beta}\) is symmetric.  This explains both the
  symmetry of \(\delta^{IJ}\) and the antisymmetry of \(Z^{IJ}\).  The
  remaining brackets may be written
  \[
  \begin{aligned}
    [P_\mu,P_\nu]&=0,&
    [P_\mu,Q^I_\alpha]&=0,\\
    [J_{\mu\nu},Q^I_\alpha]
      &=-\frac{\ii}{2}(\gamma_{\mu\nu})_\alpha{}^\beta Q^I_\beta,&
    [B_A,Q^I_\alpha]&=\ii(t_A)^I{}_JQ^J_\alpha .
  \end{aligned}
  \]
  The \(J_{\mu\nu},P_\mu\) brackets are the ordinary \(2+1\)-dimensional
  Poincar\'e algebra.  The \(B_A\) generate a compact internal algebra,
  commute with \(J_{\mu\nu}\) and \(P_\mu\), and act by antisymmetric
  matrices \(t_A\) on the real index \(I\).  The graded Jacobi identities
  require this action to preserve both \(\delta^{IJ}\) and \(Z^{IJ}\).
  The \(Z^{IJ}\) may be taken central (or, equivalently, as Abelian internal
  charges in the center of the full internal algebra).  Up to conventions
  for \(C\), \(\gamma^\mu\), and the factor of \(2\), this is the most
  general super-Poincar\'e algebra allowed by the stated hypotheses.

  \WeinbergSolution{3}
  Choose a lightlike momentum and a Clifford vacuum of helicity
  \(\lambda_{\max}\).  The \(N\) non-vanishing lowering operators each lower
  helicity by \(1/2\), so the level with \(k\) operators has
  \(\binom Nk\) states and helicity \(\lambda_{\max}-k/2\).

  For \(N=6\), the helicity range has width \(3\).  The assumed cutoff
  therefore forces \(\lambda_{\max}=3/2\), and the representation is already
  CPT self-conjugate:
  \[
  \begin{array}{c|rrrrrrr}
    \lambda&\frac32&1&\frac12&0&-\frac12&-1&-\frac32\\ \hline
    \text{multiplicity}&1&6&15&20&15&6&1 .
  \end{array}
  \]
  For \(N=5\), a multiplet beginning at \(3/2\) ends at \(-1\), with
  multiplicities \(1,5,10,10,5,1\).  CPT requires its conjugate, which
  begins at \(+1\) and ends at \(-3/2\):
  \[
  \begin{array}{c|rrrrrrr}
    \lambda&\frac32&1&\frac12&0&-\frac12&-1&-\frac32\\ \hline
    \mathcal M&1&5&10&10&5&1&0\\
    \text{CPT}(\mathcal M)&0&1&5&10&10&5&1\\ \hline
    \text{total}&1&6&15&20&15&6&1 .
  \end{array}
  \]
  The last row is exactly the \(N=6\) spectrum, by Pascal's identity
  \(\binom5k+\binom5{k-1}=\binom6k\).  Thus a CPT-invariant massless
  \(N=5\) particle spectrum with this helicity bound has the particle
  content of \(N=6\); the comparison strongly suggests enhancement from
  \(N=5\) to \(N=6\) supersymmetry.

  \WeinbergSolution{4}
  At the \(N=2\) BPS bound, half of the four complex rest-frame oscillator
  components have zero norm and act trivially.  The surviving creation
  operators form a single spin-\(1/2\) doublet, so a short \(N=2\)
  representation has the same spin-and-parity construction as a massive
  simple-supersymmetry representation.  Parity can be assigned within a
  central-charge sector only after choosing the phase of the central charge
  compatibly; otherwise parity pairs the sector with the one of opposite
  parity-transformed central charge.

  For the short gauge multiplet, take the Clifford vacuum to have
  \(j=1/2\).  The spin \(1\) and spin \(0\) particles have a common
  intrinsic parity \(\eta\), while the two spin-\(1/2\) particles have
  parities
  \[
    +\ii\eta\qquad\hbox{and}\qquad-\ii\eta.
  \]
  For the short hypermultiplet, the collapsed \(j=0\) construction gives
  one spin-\(1/2\) particle of parity \(\eta\), and the two spin-zero
  particles have parities
  \[
    +\ii\eta\qquad\hbox{and}\qquad-\ii\eta.
  \]
  The overall phase \(\eta\) is conventional, subject to the usual
  spin-statistics relation for \(P^2\); the relative factors displayed
  above are fixed by the supersymmetry algebra.
\end{enumerate}
